\apartado{Riesgos, conclusiones y trabajo pendiente}
\small
\begin{tabularx}{\linewidth}{P{2.8cm}Y Y}
\toprule
\textbf{Riesgo} & \textbf{Consecuencia} & \textbf{Tratamiento previsto}\\
\midrule
Autorización incompleta & Lectura o modificación de información ajena. & Actor del servidor, alcance por propietario y pruebas A/B/C.\\
Runtime excesivamente permisivo & Código generado alcanza datos o capacidades del navegador. & VM limitada, origen separado, SDK y ataques antes de integrar.\\
Extracción imprecisa & Fórmulas o citas erróneas sustentan explicaciones. & Corpus, localizadores, revisión y versiones de transcripción.\\
Reintentos sin control & Trabajos duplicados o resultados inconsistentes. & Idempotencia y estado persistente; integración posterior obligatoria.\\
Consumo inesperado & Excedentes cloud o llamadas IA costosas. & Cuotas, límites, medición y capacidad de detener nuevas generaciones.\\
Confundir demo y producto & Se atribuye calidad o despliegue sin evidencia. & Estado por tarea, pruebas identificadas y límites explícitos.\\
Sobrecarga de interfaz & La decoración dificulta lectura y uso móvil. & Jerarquía, teclado, contraste y evaluación de tareas reales.\\
\bottomrule
\end{tabularx}
\normalsize
\subsection{Conclusiones del diseño}
La arquitectura propuesta relaciona necesidades académicas con capacidades de nube: identidad para separar cuentas, persistencia para recuperar organización, contenedores para empaquetar ejecución y servicios diferenciados para aislar cargas. La asincronía y la observabilidad tienen función concreta en el futuro procesamiento documental; no se añaden solo como nombres de productos.

H0 reduce incertidumbre antes de ampliar funcionalidad; H1 aporta un espacio académico útil aun sin tutor. Esa secuencia permite evaluar valor y riesgos por separado. La propuesta de valor combina organización y aprendizaje conectado, pero sus beneficios siguen siendo hipótesis. No se concluye que el producto mejore calificaciones ni que supere a soluciones existentes.

La selección de servicios administrados reduce tareas de infraestructura, mientras el equipo conserva responsabilidad sobre permisos, modelos, validación y costos. El Free Tier es una restricción de planificación, no garantía de costo cero. Un presupuesto solo de alertas tampoco sustituye límites técnicos ni seguimiento de consumo.

\subsection{Trabajo pendiente}
El cierre efectivo depende del estado y evidencias del apartado final. Tras H0/H1 siguen carga y tutoría, recursos, evaluación y progreso, privacidad social, borrado completo e integración cloud verificable. Antes de un piloto con información valiosa se debe acordar retención, recuperación y uso de datos con el proveedor de IA. Quedan por decidir el nombre y el catálogo definitivo de estampillas; esos elementos no justifican inventar resultados o retrasar la validación técnica.
